\section{Literature Review}
\label{sec:literature}

What constitutes an observation in a financial news corpus depends in
part on how the news was collected. In constructing FNSPID,
\citet{dong2024} begin with the headlines and URLs on Nasdaq's
per-symbol pages, retrieve bodies from the linked pages, and supplement
the collection with earlier news data. Their modeling experiments start
from 15.7 million time-aligned news records, for which they separately
examine the availability of text and stock symbols. Collecting by symbol
gives the researcher a route from the security to the story. Starting
from a web archive reverses that order: the researcher observes a copy
of the text and must establish which, if any, securities it concerns.
A story indexed under two symbols can supply two company observations,
while two captures of the same page can document a revision. Counts of
time-aligned records and archived copies consequently measure different
aspects of news coverage. They also suggest different checks on
completeness. Starting with a list of securities, one can ask which
ones have any retrievable text at all. Starting with archived articles,
one can ask how often a recognized company resolves to a listed issuer.
To learn which firms receive little or no coverage, however, both
approaches need a roster defined independently of the articles found.
Otherwise an issuer with no stories disappears from the denominator
along with the evidence of its absence.

For researchers starting with the archived page, the choice of corpus
also determines which editorial and processing decisions they inherit.
The daily CC-News files preserve fetched responses from news sites
\citep{nagel2016}, while Infini-News extracts article text and metadata
that can be searched together with their capture provenance
\citep{lazzaroni2026}. Fine News further filters the extracted records by
language and quality and selects representatives of duplicate texts
\citep{solovev2026}. For pretraining, filtering and deduplication help
determine which passages a model encounters and how often it encounters
them \citep{penedo2023}. For a financial panel, the same operations
can change the measured volume of coverage. Two publishers carrying the
same wire story may each contribute a publication to a study of media
attention, although counting both gives an additional observation of the
same bank pair in a study of co-mentions. The appropriate treatment of
the second copy depends on which of these quantities the study estimates.

A further distinction arises between finding an organization in the
text and identifying a security that a researcher can follow over time.
GDELT's Global Knowledge Graph records organizations, themes and tone
in monitored news and therefore allows researchers to follow how events
are reported \citep{gdelt}. FiNER-ORD makes recognition itself amenable
to evaluation: \citet{shah2024} provide manual annotations for people,
organizations and locations in 201 financial news articles. Those labels
let a researcher ask whether an extractor found the right words in a
passage. In a financial application, a correctly extracted organization
might be an exchange, a regulator or a subsidiary with a listed parent,
each of which has a different implication for whether the mention can
enter an issuer panel. \citet{schwenkler2020} make this additional step
explicit when they recognize organizations in news and link firms to
construct a network. A missed name can remove a genuine pair, while an
incorrect assignment can attribute that pair to a different firm.

Even a correct assignment leaves open what the text says about the firm.
Of the 10,753 human-annotated headlines in SEntFiN, 2,847 contain several
entities, and the sentiments expressed toward them can differ
\citep{sinha2022sentfin}. Consider a headline announcing the acquisition
of one company by another. The buyer and the target should both be
retrievable by name, though an investor interpreting the terms of the
transaction may draw different conclusions for each security.
\citet{heston2017} can study such firm-level questions by matching news
text to vendor firm tags and sentiment scores. For archived text, the
corresponding decisions can be examined separately by asking whether
the name was linked to the right issuer, whether the story concerned
that issuer, and whether the sentiment assigned to it agrees with a
judgment made in context.

The financial application determines which of these decisions enters
the estimate. In the volatility tracker of \citet{baker2026}, articles
that meet a dictionary rule contribute to a publisher's monthly share,
so the other articles the publisher carried that month also enter the
calculation through the denominator. In the bank networks of
\citet{kazinnik2026}, a connection requires evidence about two
institutions in the same story. The overnight return design of
\citet{lopez2026} adds a time requirement, since a headline can inform
a trading decision only if it was available before the position was
taken. For a captured article body, that requirement concerns the
version containing the information used by the model. We therefore
begin the construction of CCFinance with the archived copy and carry
its identity forward as we assign issuers, evaluate their roles in
the story and establish which observations can enter each design.
